\section{Limitations}
\label{sec:limits}

This study covers two smooth Hamiltonian systems, a pendulum and a two-degree-of-freedom anharmonic
oscillator, and neither has contacts or multiple interacting objects, nor have we tested
discontinuous dynamics. Each arm has three trained models, and the model seed is the independent
unit. Three is too few for a precise effect size, so we report ranges rather than confidence
intervals. The extraction minimises an invariance ratio that \emph{is} the normal component of
one-step error, so a low normal error shows the optimiser meeting its objective, not a finding. The
empirical content is that the resulting direction coincides with physical energy, which supervision
does not reproduce. The architecture comparison likewise does not isolate architecture, because an
open-loop reconstruction term and a prior--posterior KL are different losses, not the same loss at
different strengths.

\paragraph{No downstream benefit, and we measured why.}
Both practical uses we tested fail, for one reason. We plan with CEM over the learned model on an
energy-targeting task, and the model plans well, beating random actions by $+1.65$, 95\% CI
$[+0.50,+2.79]$. Yet no correction arm differs from no correction over three models and 20 episodes
each, and the largest effect is $0.0014$ against a return standard deviation of $0.575$, or
\textbf{0.24\% of one episode's spread}. All four registered predictions fail. The cause is
structural: the correction acts on \emph{accumulated} drift, and at the planner's horizon it shifts
imagined energy by $0.3\%$ of the spread across candidate action sequences. The planner ranks plans
by that spread, so a shift this small cannot change which plan wins. We measure our repair over 100
autonomous steps, and a re-planning controller never imagines long enough to accumulate what the
method removes. Looking further ahead does not help, because planning quality itself degrades with
horizon. The same effect-size limit defeats drift as an online trust signal, which plain latent
displacement beats on every model, and Appendix~\ref{sec:causal-detail} gives both in full.

\paragraph{The method finds only what the data gave the model a reason to represent.}
The pendulum's rod length never changes, an exact constraint, yet a search for a residual $G(z)=0$
fails, and the recovered direction correlates with the true residual at $0.004$--$0.032$. Minimising
$\mathbb{E}[G^2]$ is the wrong objective, ranking the true constraint $1525$th of $1819$ directions.
The model barely encodes the constraint at all, and a probe predicts it from the latent at held-out
$|\rho|=0.19$. \textbf{A quantity that never varies carries no information}, so nothing in a
reconstruction objective pushes it into the latent, and it can live in the decoder's weights. This
bounds the approach to quantities that \emph{vary}, and it is why we report no released-checkpoint
experiment: limb lengths do not vary either.

Finally, the two-invariant case limits the extraction rather than the model. Where the system
conserves both energy and angular momentum, the conserved subspace reliably contains both, at
angular-momentum correlation $0.94$--$0.98$, while the joint fit isolates either in only 1 of 7
measurements. Flow alignment selects a single direction inside that subspace, and it has no unique
answer when two exist.

\paragraph{We cannot attribute the shadow correction to the model rather than the data.}
Section~\ref{sec:mechanism} establishes the account on the data side and reports, in the same place,
why the model side resists measurement here. We compute $\rho_{\mathrm{obs}}$ along evaluation
trajectories, and a registered cross-evaluation control shows the argmin follows the evaluation set
in $6/6$ cells and the checkpoint in $0/6$. Of the two alternatives that remain, the rollout
statistic has no power, at a contrast of $1.12$--$1.45\times$. The cross-timestep forced choice is
unreadable, because a model evaluated off its training timestep is $\sim\!90\times$ worse in one-step
error. Deciding the question needs a model in distribution at more than one timestep. Until then the
claim is unsubstantiated rather than refuted, and we construct the paper's other results so that they
do not depend on it.
